% problem_id: S001-lemma-blowup
% source_id: S001 (Stacks Project)
% source_file: algebraization.tex
% statement_locator: algebraization.tex lines 7292--7315
% proof_locator: algebraization.tex lines 7317--7581
% env: lemma
% pairing: tight (verified)
% status: complete (statement + author proof verbatim + reason + locators)
%
\begin{lemma}
\label{lemma-blowup}
In Situation \ref{situation-algebraize} let $(\mathcal{F}_n)$
be an object of $\textit{Coh}(U, I\mathcal{O}_U)$. Assume
\begin{enumerate}
\item $A$ has a dualizing complex,
\item all fibres of the blowing up $b : X' \to X$ of $I$
have dimension $\leq d - 1$,
\item one of the following is true
\begin{enumerate}
\item $(\mathcal{F}_n)$ satisfies the $(d + 1, d + 2)$-inequalities
(Definition \ref{definition-s-d-inequalities}), or
\item for $y \in U \cap Y$ and a prime
$\mathfrak p \subset \mathcal{O}_{X, y}^\wedge$ with
$\mathfrak p \not \in V(I\mathcal{O}_{X, y}^\wedge)$
we have
$$
\text{depth}((\mathcal{F}^\wedge_y)_\mathfrak p) +
\dim(\mathcal{O}_{X, y}^\wedge/\mathfrak p) + \delta^Y_Z(y) > d + 2
$$
\end{enumerate}
\end{enumerate}
Then $(\mathcal{F}_n)$ extends to $X$.
\end{lemma}

\begin{proof}
Let $Y' \subset X'$ be the exceptional divisor.
Let $Z' \subset Y'$ be the inverse image of $Z \subset Y$.
Then $U' = X' \setminus Z'$ is the inverse image of $U$.
With $\delta^{Y'}_{Z'}$ as in (\ref{equation-delta-Z}) we set
$$
T' = \{y' \in Y' \mid \delta^{Y'}_{Z'}(y') = 1\}
\subset
T = \{y' \in Y' \mid \delta^{Y'}_{Z'}(y') = 1\text{ or }2\}
$$
By Lemma \ref{lemma-discussion} parts (1) and (2) these are specialization
stable subsets of $U' \cap Y' = Y' \setminus Z'$. Consider the
object $(b|_{U'}^*\mathcal{F}_n)$ of $\textit{Coh}(U', I\mathcal{O}_{U'})$,
see Cohomology of Schemes, Lemma \ref{coherent-lemma-inverse-systems-pullback}.
For $y' \in U' \cap Y'$ let us denote
$$
\mathcal{F}_{y'}^\wedge = \lim (b|_{U'}^*\mathcal{F}_n)_{y'}
$$
the ``stalk'' of this pullback at $y'$. We claim that conditions
(a), (b), (c), (d), and (e) of
Lemma \ref{lemma-improvement-formal-coherent-module-better}
hold for the object $(b|_{U'}^*\mathcal{F}_n)$ on $U'$ with $d$
replaced by $1$ and the subsets $T' \subset T \subset U' \cap Y'$.
Condition (a) holds because $Y'$ is an effective Cartier divisor
and hence locally cut out by $1$ equation. Condition (e) holds
by Lemma \ref{lemma-discussion} part (7).
To prove (b), (c), and (d) we need some preparation.

\medskip\noindent
Let $y' \in U' \cap Y'$ and let
$\mathfrak p' \subset \mathcal{O}_{X', y'}^\wedge$
be a prime ideal not contained in $V(I\mathcal{O}_{X', y'}^\wedge)$.
Denote $y = b(y') \in U \cap Y$. Choose $f \in I$ such that
$y'$ is contained in the spectrum of the affine blowup algebra
$A[\frac{I}{f}]$, see Divisors, Lemma \ref{divisors-lemma-blowing-up-affine}.
For any $A$-algebra $B$ denote $B' = B[\frac{IB}{f}]$ the corresponding affine
blowup algebra. Denote $I$-adic completion by ${\ }^\wedge$.
By our choice of $f$ we get a ring map
$(\mathcal{O}_{X, y}^\wedge)' \to \mathcal{O}_{X', y'}^\wedge$.
If we let $\mathfrak q' \subset (\mathcal{O}_{X, y}^\wedge)'$
be the inverse image of $\mathfrak m_{y'}^\wedge$, then
we see that
$((\mathcal{O}_{X, y}^\wedge)'_{\mathfrak q'})^\wedge =
\mathcal{O}_{X', y'}^\wedge$.
Let $\mathfrak p \subset \mathcal{O}_{X, y}^\wedge$ be the corresponding
prime. At this point we have a commutative diagram
$$
\xymatrix{
\mathcal{O}_{X, y}^\wedge \ar[d] \ar[r] &
(\mathcal{O}_{X, y}^\wedge)' \ar[d]_\alpha \ar[r] &
(\mathcal{O}_{X, y}^\wedge)'_{\mathfrak q'} \ar[d] \ar[r]_\beta &
\mathcal{O}_{X', y'}^\wedge \ar[d] \\
\mathcal{O}_{X, y}^\wedge/\mathfrak p \ar[r] &
(\mathcal{O}_{X, y}^\wedge/\mathfrak p)' \ar[r] &
(\mathcal{O}_{X, y}^\wedge/\mathfrak p)'_{\mathfrak q'} \ar[r]^\gamma &
((\mathcal{O}_{X, y}^\wedge/\mathfrak p)'_{\mathfrak q'})^\wedge \ar[d] \\
 & & &
\mathcal{O}_{X', y'}^\wedge/\mathfrak p'
}
$$
whose vertical arrows are surjective. By
More on Algebra, Lemma \ref{more-algebra-lemma-completion-dimension}
and the dimension formula
(Algebra, Lemma \ref{algebra-lemma-dimension-formula})
we have
$$
\dim(((\mathcal{O}_{X, y}^\wedge/\mathfrak p)'_{\mathfrak q'})^\wedge) =
\dim((\mathcal{O}_{X, y}^\wedge/\mathfrak p)'_{\mathfrak q'}) =
\dim(\mathcal{O}_{X, y}^\wedge/\mathfrak p)
- \text{trdeg}(\kappa(y')/\kappa(y))
$$
Tracing through the definitions of pullbacks, stalks, localizations,
and completions we find
$$
(\mathcal{F}_y^\wedge)_{\mathfrak p}
\otimes_{(\mathcal{O}_{X, y}^\wedge)_\mathfrak p}
(\mathcal{O}_{X', y'}^\wedge)_{\mathfrak p'}
=
(\mathcal{F}_{y'}^\wedge)_{\mathfrak p'}
$$
Details omitted. The ring maps $\beta$ and $\gamma$ in the diagram
are flat with Gorenstein (hence Cohen-Macaulay) fibres, as these are
completions of rings having a dualizing complex. See
Dualizing Complexes, Lemmas
\ref{dualizing-lemma-formal-fibres-gorenstein} and
\ref{dualizing-lemma-dualizing-gorenstein-formal-fibres}
and the discussion in More on Algebra, Section
\ref{more-algebra-section-properties-formal-fibres}.
Observe that $(\mathcal{O}_{X, y}^\wedge)_\mathfrak p =
(\mathcal{O}_{X, y}^\wedge)'_{\tilde{\mathfrak p}}$
where $\tilde{\mathfrak p}$ is the kernel of $\alpha$
in the diagram. On the other hand,
$(\mathcal{O}_{X, y}^\wedge)'_{\tilde{\mathfrak p}}
\to (\mathcal{O}_{X', y'}^\wedge)_{\mathfrak p'}$
is flat with CM fibres by the above. Whence
$(\mathcal{O}_{X, y}^\wedge)_\mathfrak p  \to
(\mathcal{O}_{X', y'}^\wedge)_{\mathfrak p'}$ is flat with CM fibres.
Using Algebra, Lemma \ref{algebra-lemma-apply-grothendieck-module}
we see that
$$
\text{depth}((\mathcal{F}_{y'}^\wedge)_{\mathfrak p'}) =
\text{depth}((\mathcal{F}_y^\wedge)_{\mathfrak p}) +
\dim(F_\mathfrak r)
$$
where $F$ is the generic formal fibre of
$(\mathcal{O}_{X, y}^\wedge/\mathfrak p)'_{\mathfrak q'}$
and $\mathfrak r$ is the prime corresponding to $\mathfrak p'$.
Since $(\mathcal{O}_{X, y}^\wedge/\mathfrak p)'_{\mathfrak q'}$
is a universally catenary local domain, its $I$-adic completion
is equidimensional and (universally) catenary by Ratliff's theorem
(More on Algebra, Proposition \ref{more-algebra-proposition-ratliff}).
It then follows that
$$
\dim(((\mathcal{O}_{X, y}^\wedge/\mathfrak p)'_{\mathfrak q'})^\wedge) =
\dim(F_\mathfrak r) + \dim(\mathcal{O}_{X', y'}^\wedge/\mathfrak p')
$$
Combined with Lemma \ref{lemma-change-distance-function}
we get
\begin{equation}
\label{equation-one}
\begin{aligned}
&
\text{depth}((\mathcal{F}_{y'}^\wedge)_{\mathfrak p'}) +
\delta^{Y'}_{Z'}(y') \\
& =
\text{depth}((\mathcal{F}_y^\wedge)_{\mathfrak p}) +
\dim(F_\mathfrak r) + \delta^{Y'}_{Z'}(y') \\
& \geq
\text{depth}((\mathcal{F}_y^\wedge)_{\mathfrak p}) + \delta^Y_Z(y) +
\dim(F_\mathfrak r) + \text{trdeg}(\kappa(y')/\kappa(y)) - (d - 1) \\
& =
\text{depth}((\mathcal{F}_y^\wedge)_{\mathfrak p}) + \delta^Y_Z(y) - (d - 1)
+ \dim(\mathcal{O}_{X, y}^\wedge/\mathfrak p) -
\dim(\mathcal{O}_{X', y'}^\wedge/\mathfrak p')
\end{aligned}
\end{equation}
Please keep in mind that
$\dim(\mathcal{O}_{X, y}^\wedge/\mathfrak p) \geq
\dim(\mathcal{O}_{X', y'}^\wedge/\mathfrak p')$. Rewriting this we get
\begin{equation}
\label{equation-two}
\begin{aligned}
&
\text{depth}((\mathcal{F}_{y'}^\wedge)_{\mathfrak p'}) +
\dim(\mathcal{O}_{X', y'}^\wedge/\mathfrak p') +
\delta^{Y'}_{Z'}(y') \\
& \geq
\text{depth}((\mathcal{F}_y^\wedge)_{\mathfrak p}) +
\dim(\mathcal{O}_{X, y}^\wedge/\mathfrak p) +
\delta^Y_Z(y) - (d - 1)
\end{aligned}
\end{equation}
This inequality will allow us to check the remaining conditions.

\medskip\noindent
Conditions (b) and (d) of
Lemma \ref{lemma-improvement-formal-coherent-module-better}. Assume
$V(\mathfrak p') \cap V(I\mathcal{O}_{X', y'}^\wedge) =
\{\mathfrak m_{y'}^\wedge\}$.
This implies that $\dim(\mathcal{O}_{X', y'}^\wedge/\mathfrak p') = 1$
because $Z'$ is an effective Cartier divisor.
The combination of (b) and (d) is equivalent with
$$
\text{depth}((\mathcal{F}_{y'}^\wedge)_{\mathfrak p'}) + \delta^{Y'}_{Z'}(y')
> 2
$$
If $(\mathcal{F}_n)$ satisfies the inequalities in (3)(b)
then we immediately conclude this is true by applying (\ref{equation-two}).
If $(\mathcal{F}_n)$ satisfies (3)(a), i.e., the
$(d + 1, d + 2)$-inequalities, then we see that in any case
$$
\text{depth}((\mathcal{F}_y^\wedge)_{\mathfrak p}) + \delta^Y_Z(y)
\geq d + 1
\quad\text{or}\quad
\text{depth}((\mathcal{F}_y^\wedge)_{\mathfrak p}) +
\dim(\mathcal{O}_{X, y}^\wedge/\mathfrak p) +
\delta^Y_Z(y) > d + 2
$$
Looking at (\ref{equation-one}) and (\ref{equation-two}) above this gives
what we want except possibly if
$\dim(\mathcal{O}_{X, y}^\wedge/\mathfrak p) = 1$.
However, if $\dim(\mathcal{O}_{X, y}^\wedge/\mathfrak p) = 1$, then we have
$V(\mathfrak p) \cap V(I\mathcal{O}_{X, y}^\wedge) = \{\mathfrak m_y^\wedge\}$
and we see that actually
$$
\text{depth}((\mathcal{F}_y^\wedge)_{\mathfrak p}) + \delta^Y_Z(y) > d + 1
$$
as $(\mathcal{F}_n)$ satisfies the $(d + 1, d + 2)$-inequalities and we
conclude again.

\medskip\noindent
Condition (c) of
Lemma \ref{lemma-improvement-formal-coherent-module-better}. Assume
$V(\mathfrak p') \cap V(I\mathcal{O}_{X', y'}^\wedge) \not =
\{\mathfrak m_{y'}^\wedge\}$. Then condition (c) is equivalent to
$$
\text{depth}((\mathcal{F}_{y'}^\wedge)_{\mathfrak p'}) + \delta^{Y'}_{Z'}(y')
\geq 2
\quad\text{or}\quad
\text{depth}((\mathcal{F}_{y'}^\wedge)_{\mathfrak p'}) +
\dim(\mathcal{O}_{X', y'}^\wedge/\mathfrak p') +
\delta^{Y'}_{Z'}(y') > 3
$$
If $(\mathcal{F}_n)$ satisfies the inequalities in (3)(b)
then we see the second of the two displayed inequalities holds true
by applying (\ref{equation-two}). If $(\mathcal{F}_n)$ satisfies (3)(a), i.e.,
the $(d + 1, d + 2)$-inequalities, then this follows immediately from
(\ref{equation-one}) and (\ref{equation-two}).
This finishes the proof of our claim.

\medskip\noindent
Choose $(b|_{U'}^*\mathcal{F}_n) \to (\mathcal{F}_n'')$
and $(\mathcal{H}_n)$ in $\textit{Coh}(U', I\mathcal{O}_{U'})$
as in Lemma \ref{lemma-improvement-formal-coherent-module-better}.
For any affine open $W \subset X'$ observe that
$\delta^{W \cap Y'}_{W \cap Z'}(y') \geq \delta^{Y'}_{Z'}(y')$ by
Lemma \ref{lemma-discussion} part (6). Hence we see that
$(\mathcal{H}_n|_W)$ satisfies the assumptions of
Lemma \ref{lemma-cd-1-canonical}.
Thus $(\mathcal{H}_n|_W)$ extends canonically to $W$.
Let $(\mathcal{G}_{W, n})$ in $\textit{Coh}(W, I\mathcal{O}_W)$
be the canonical extension as in
Lemma \ref{lemma-canonically-algebraizable}.
By Lemma \ref{lemma-canonically-extend-base-change}
we see that for $W' \subset W$ there is a unique isomorphism
$$
(\mathcal{G}_{W, n}|_{W'}) \longrightarrow
(\mathcal{G}_{W', n})
$$
compatible with the given isomorphisms
$(\mathcal{G}_{W, n}|_{W \cap U}) \cong (\mathcal{H}_n|_{W \cap U})$.
We conclude that there exists an object
$(\mathcal{G}_n)$ of $\textit{Coh}(X', I\mathcal{O}_{X'})$
whose restriction to $U$ is isomorphic to $(\mathcal{H}_n)$.

\medskip\noindent
If $A$ is $I$-adically complete we can finish the proof as follows.
By Grothendieck's existence theorem
(Cohomology of Schemes, Lemma \ref{coherent-lemma-existence-projective})
we see that $(\mathcal{G}_n)$ is the completion of a coherent
$\mathcal{O}_{X'}$-module. Then by
Cohomology of Schemes, Lemma \ref{coherent-lemma-existence-easy}
we see that $(b|_{U'}^*\mathcal{F}_n)$
is the completion of a coherent $\mathcal{O}_{U'}$-module
$\mathcal{F}'$.
By Cohomology of Schemes, Lemma \ref{coherent-lemma-inverse-systems-push-pull}
we see that there is a map
$$
(\mathcal{F}_n) \longrightarrow ((b|_{U'})_*\mathcal{F}')^\wedge
$$
whose kernel and cokernel is annihilated by a power of $I$.
Then finally, we win by applying
Lemma \ref{lemma-map-kernel-cokernel-on-closed}.

\medskip\noindent
If $A$ is not complete, then, before starting the proof, we may replace $A$
by its completion, see Lemma \ref{lemma-algebraizable}.
After completion the assumptions still hold: this is immediate
for condition (3), follows from
Dualizing Complexes, Lemma \ref{dualizing-lemma-ubiquity-dualizing}
for condition (1), and from
Divisors, Lemma \ref{divisors-lemma-flat-base-change-blowing-up}
for condition (2).
Thus the complete case implies the general case.
\end{proof}
